\chapter{Problem Statement}
\label{ch:problem}

\section{The Existing Workflow and its Pain Points}
\label{sec:existing-workflow}

Consider a typical academic scenario in which a student receives a
tabular dataset (e.g.\ a Kaggle CSV, a lab measurement export, or a
survey response file) with the instruction ``build a predictive model
and analyse it''. The student's usual workflow involves:

\begin{enumerate}[leftmargin=1.4em]
  \item Opening a fresh Jupyter notebook.
  \item Importing pandas, scikit-learn, seaborn or matplotlib, and
        possibly XGBoost / LightGBM.
  \item Manually computing \code{df.describe()}, \code{df.isna()},
        \code{df.corr()} and plotting a handful of charts.
  \item Splitting the data with \code{train\_test\_split}, choosing a
        model almost arbitrarily, fitting it, and reporting one or two
        metrics.
  \item Rarely: computing feature importances or SHAP values.
  \item Almost never: producing a self-contained, executable notebook
        or PDF report that can be graded, reviewed, or reproduced.
\end{enumerate}

The failure modes are well documented in the literature
\citep{pineau2021reproducibility, kapoor2022leakage} and observed
first-hand in classroom settings: incorrect train/test splits, target
leakage, inappropriate metrics, hidden preprocessing, unstated random
seeds, and missing environment specifications. These failures are
rarely due to bad intent; they arise from cognitive load and from the
absence of an opinionated, guided workflow.

\section{Target Users}
\label{sec:target-users}

\projname{} is designed for three overlapping audiences:

\begin{itemize}[leftmargin=1.4em]
  \item \textbf{Students} (undergraduate and postgraduate) doing
        coursework or capstone projects with tabular data.
  \item \textbf{Applied researchers} (in economics, biology, social
        sciences, engineering) who need a quick baseline before
        committing to a bespoke model.
  \item \textbf{Data engineers} exploring MCP as a mechanism for
        exposing internal data products to LLM agents.
\end{itemize}

The platform is explicitly \emph{not} designed for large-scale
production use, for competitive Kaggle-style benchmarking, or for
non-tabular data.

\section{Problem Statement}
\label{sec:formal-problem}

We can now state the problem this project addresses:

\begin{quote}
\emph{Design and implement a local, single-user, extensible platform
that lets a non-expert user upload a tabular dataset, ask
natural-language questions and issue natural-language commands, and
receive back --- with an audit trail of tool calls, arguments, and
intermediate results --- a full set of exploratory analyses, a
comparative model leaderboard, an explainability card, a reproducible
Jupyter notebook, and a PDF report, all persisted on the local
filesystem.}
\end{quote}

Concretely, the platform must satisfy the following relationships:

\begin{itemize}[leftmargin=1.4em]
  \item Given a natural-language request $q$ and an active dataset
        $D$, the platform selects a sequence of tools
        $t_1,\ldots,t_k$ from a fixed catalog and produces a response
        $r(q,D) = \Phi(t_k \circ \cdots \circ t_1)(D)$ together with
        an audit trace $[(t_i, a_i, o_i)]$ where $a_i$ are the
        arguments and $o_i$ the outputs.
  \item Every artefact produced by any $t_i$ is written to the local
        filesystem with a UUID and metadata sufficient to recover it
        later.
  \item The system must remain usable when the LLM is unavailable ---
        the tool catalog can be invoked programmatically without a
        chat turn.
\end{itemize}

\section{Design Goals}
\label{sec:design-goals}

The design of \projname{} is guided by the following goals, in order
of priority:

\begin{enumerate}[leftmargin=1.4em]
  \item \textbf{Determinism first.} The LLM plans and formats; every
        numerical result comes from deterministic Python code with a
        fixed random seed.
  \item \textbf{Auditability.} Every tool call is visible to the user
        (arguments, result, duration, service).
  \item \textbf{Reproducibility.} Every generated artefact is stored
        with a stable identifier and metadata.
  \item \textbf{Extensibility.} Adding a new tool is a
        single-decorator change to one MCP server; nothing else needs
        to change.
  \item \textbf{Local-first.} No cloud storage, no authentication, no
        network I/O except the LLM API.
  \item \textbf{Thin UI.} The Streamlit layer contains no ML logic;
        it is a client of the orchestrator.
\end{enumerate}

\section{Constraints}
\label{sec:constraints}

The following constraints were adopted throughout the design:

\begin{itemize}[leftmargin=1.4em]
  \item \textbf{Language.} All server-side code is in Python
        3.11+. This aligns with the modern typing features used by
        FastMCP.
  \item \textbf{Runtime.} Single machine, single OS process.
        Concurrency is exclusively via \code{asyncio} and small
        thread pools used inside individual tools.
  \item \textbf{Storage.} Local filesystem under \file{data/} only.
        No database.
  \item \textbf{LLM provider.} Groq's OpenAI-compatible Chat
        Completions API, primarily with Llama\,3.3 70B; other
        OpenAI-compatible providers are drop-in but untested.
  \item \textbf{Data modality.} Tabular only. Image, text, audio,
        and multimodal data are out of scope.
\end{itemize}

\section{Assumptions}
\label{sec:assumptions}

The design assumes that:

\begin{itemize}[leftmargin=1.4em]
  \item The user has a working Python 3.11+ environment and can
        install the listed dependencies with \code{pip}.
  \item The user has a Groq API key. The system explicitly warns
        the user when the key is missing.
  \item The dataset fits comfortably in memory (default upload cap
        is 200~MB, configurable).
  \item The user is trusted --- uploaded files may execute pandas
        parsing but not arbitrary code.
\end{itemize}

\section{Success Criteria}
\label{sec:success}

The platform is considered successful for the purposes of this project
if, on a standard tabular dataset (e.g.\ the UCI Iris
\citep{uciiris}, Wine \citep{uciwine}, or Adult \citep{uciadult}
datasets), a first-time user can:

\begin{enumerate}[leftmargin=1.4em]
  \item Upload the CSV and see it appear in the sidebar.
  \item Issue a single natural-language request such as ``run a full
        EDA and then train a model to predict \code{<target>}''.
  \item Observe the copilot select the correct sequence of tools,
        display live progress, and produce (a)~an EDA report with
        three plots, (b)~a champion model, (c)~a SHAP summary plot
        or feature-importance chart, (d)~a Jupyter notebook that
        reproduces the pipeline, and (e)~a PDF report that lists the
        leaderboard.
  \item Download each artefact from the artefacts tab.
\end{enumerate}

The remainder of the report shows how these criteria are met.
